% Figure 4: Ablation Study - Iteration vs Performance
% 사용법: % Figure 4: Ablation Study - Iteration vs Performance
% 사용법: % Figure 4: Ablation Study - Iteration vs Performance
% 사용법: % Figure 4: Ablation Study - Iteration vs Performance
% 사용법: \input{figures/fig_ablation.tex}
\begin{figure}[t]
\centering
\resizebox{\columnwidth}{!}{%
\begin{tikzpicture}
\begin{axis}[
    width=10cm,
    height=6.5cm,
    xlabel={반복 횟수 (Iterations)},
    ylabel={해결 문제 수 (/100)},
    xlabel style={font=\small\sffamily},
    ylabel style={font=\small\sffamily},
    xticklabel style={font=\small\sffamily},
    yticklabel style={font=\small\sffamily},
    xtick={1,3,5,10},
    ytick={0,5,10,15,19,20,21,22,25,28,30},
    ymin=0, ymax=32,
    xmin=0, xmax=11.5,
    grid=major,
    grid style={dashed, gray!25},
    legend style={
        at={(0.03,0.97)},
        anchor=north west,
        font=\small\sffamily,
        draw=gray!40,
        fill=white,
        fill opacity=0.9,
    },
]

% 개별 실행 성능
\addplot[
    name path=individual,
    color=blue!70,
    mark=*,
    mark size=3pt,
    thick,
    line join=round,
] coordinates {
    (1, 19) (3, 21) (5, 19) (10, 22)
};
\addlegendentry{개별 실행 성능}

% Union (누적 합집합) - 점선으로 표시
\addplot[
    name path=union,
    color=red!70,
    mark=triangle*,
    mark size=3.5pt,
    thick,
    dashed,
] coordinates {
    (1, 19) (3, 25) (5, 26) (10, 28)
};
\addlegendentry{누적 합집합 (Union)}

% 다양성 격차(shaded region)
\addplot[fill=green!10, draw=none, forget plot]
    fill between[of=union and individual];

% 개별 최대값 수평선
\addplot[
    color=blue!30,
    dashed,
    thin,
    domain=0:11.5,
] {22};

% Union 최대값 수평선
\addplot[
    color=red!30,
    dashed,
    thin,
    domain=0:11.5,
] {28};

% 차이 주석
\draw[{Stealth[length=4pt]}-{Stealth[length=4pt]}, thick, color=green!50!black]
    (10.8, 22) -- (10.8, 28);
\node[font=\scriptsize\bfseries\sffamily, color=green!50!black, right] at (10.8, 25) {+6};

% 라벨
\node[font=\tiny\sffamily, color=blue!50, anchor=west] at (11, 22) {max=22};
\node[font=\tiny\sffamily, color=red!50, anchor=west] at (11, 28) {union=28};

\end{axis}

% 하단 설명
\node[font=\scriptsize\sffamily, color=black!60, align=center] at (4.5, -0.8) {
    확률적 다양성 원리: 합집합(28) $>$ 개별 최대(22), 다양성 격차 = +27.3\%
};

\end{tikzpicture}
}
\caption{반복 횟수에 따른 성능 변화 (gpt-4o, Level~1). 개별 실행 성능은 반복 횟수에 따라 미미한 변동을 보이나, 누적 합집합은 꾸준히 증가한다. 이는 반복의 깊이보다 생성의 다양성이 더 효과적임을 시사한다.}
\label{fig:ablation}
\end{figure}

\begin{figure}[t]
\centering
\resizebox{\columnwidth}{!}{%
\begin{tikzpicture}
\begin{axis}[
    width=10cm,
    height=6.5cm,
    xlabel={반복 횟수 (Iterations)},
    ylabel={해결 문제 수 (/100)},
    xlabel style={font=\small\sffamily},
    ylabel style={font=\small\sffamily},
    xticklabel style={font=\small\sffamily},
    yticklabel style={font=\small\sffamily},
    xtick={1,3,5,10},
    ytick={0,5,10,15,19,20,21,22,25,28,30},
    ymin=0, ymax=32,
    xmin=0, xmax=11.5,
    grid=major,
    grid style={dashed, gray!25},
    legend style={
        at={(0.03,0.97)},
        anchor=north west,
        font=\small\sffamily,
        draw=gray!40,
        fill=white,
        fill opacity=0.9,
    },
]

% 개별 실행 성능
\addplot[
    name path=individual,
    color=blue!70,
    mark=*,
    mark size=3pt,
    thick,
    line join=round,
] coordinates {
    (1, 19) (3, 21) (5, 19) (10, 22)
};
\addlegendentry{개별 실행 성능}

% Union (누적 합집합) - 점선으로 표시
\addplot[
    name path=union,
    color=red!70,
    mark=triangle*,
    mark size=3.5pt,
    thick,
    dashed,
] coordinates {
    (1, 19) (3, 25) (5, 26) (10, 28)
};
\addlegendentry{누적 합집합 (Union)}

% 다양성 격차(shaded region)
\addplot[fill=green!10, draw=none, forget plot]
    fill between[of=union and individual];

% 개별 최대값 수평선
\addplot[
    color=blue!30,
    dashed,
    thin,
    domain=0:11.5,
] {22};

% Union 최대값 수평선
\addplot[
    color=red!30,
    dashed,
    thin,
    domain=0:11.5,
] {28};

% 차이 주석
\draw[{Stealth[length=4pt]}-{Stealth[length=4pt]}, thick, color=green!50!black]
    (10.8, 22) -- (10.8, 28);
\node[font=\scriptsize\bfseries\sffamily, color=green!50!black, right] at (10.8, 25) {+6};

% 라벨
\node[font=\tiny\sffamily, color=blue!50, anchor=west] at (11, 22) {max=22};
\node[font=\tiny\sffamily, color=red!50, anchor=west] at (11, 28) {union=28};

\end{axis}

% 하단 설명
\node[font=\scriptsize\sffamily, color=black!60, align=center] at (4.5, -0.8) {
    확률적 다양성 원리: 합집합(28) $>$ 개별 최대(22), 다양성 격차 = +27.3\%
};

\end{tikzpicture}
}
\caption{반복 횟수에 따른 성능 변화 (gpt-4o, Level~1). 개별 실행 성능은 반복 횟수에 따라 미미한 변동을 보이나, 누적 합집합은 꾸준히 증가한다. 이는 반복의 깊이보다 생성의 다양성이 더 효과적임을 시사한다.}
\label{fig:ablation}
\end{figure}

\begin{figure}[t]
\centering
\resizebox{\columnwidth}{!}{%
\begin{tikzpicture}
\begin{axis}[
    width=10cm,
    height=6.5cm,
    xlabel={반복 횟수 (Iterations)},
    ylabel={해결 문제 수 (/100)},
    xlabel style={font=\small\sffamily},
    ylabel style={font=\small\sffamily},
    xticklabel style={font=\small\sffamily},
    yticklabel style={font=\small\sffamily},
    xtick={1,3,5,10},
    ytick={0,5,10,15,19,20,21,22,25,28,30},
    ymin=0, ymax=32,
    xmin=0, xmax=11.5,
    grid=major,
    grid style={dashed, gray!25},
    legend style={
        at={(0.03,0.97)},
        anchor=north west,
        font=\small\sffamily,
        draw=gray!40,
        fill=white,
        fill opacity=0.9,
    },
]

% 개별 실행 성능
\addplot[
    name path=individual,
    color=blue!70,
    mark=*,
    mark size=3pt,
    thick,
    line join=round,
] coordinates {
    (1, 19) (3, 21) (5, 19) (10, 22)
};
\addlegendentry{개별 실행 성능}

% Union (누적 합집합) - 점선으로 표시
\addplot[
    name path=union,
    color=red!70,
    mark=triangle*,
    mark size=3.5pt,
    thick,
    dashed,
] coordinates {
    (1, 19) (3, 25) (5, 26) (10, 28)
};
\addlegendentry{누적 합집합 (Union)}

% 다양성 격차(shaded region)
\addplot[fill=green!10, draw=none, forget plot]
    fill between[of=union and individual];

% 개별 최대값 수평선
\addplot[
    color=blue!30,
    dashed,
    thin,
    domain=0:11.5,
] {22};

% Union 최대값 수평선
\addplot[
    color=red!30,
    dashed,
    thin,
    domain=0:11.5,
] {28};

% 차이 주석
\draw[{Stealth[length=4pt]}-{Stealth[length=4pt]}, thick, color=green!50!black]
    (10.8, 22) -- (10.8, 28);
\node[font=\scriptsize\bfseries\sffamily, color=green!50!black, right] at (10.8, 25) {+6};

% 라벨
\node[font=\tiny\sffamily, color=blue!50, anchor=west] at (11, 22) {max=22};
\node[font=\tiny\sffamily, color=red!50, anchor=west] at (11, 28) {union=28};

\end{axis}

% 하단 설명
\node[font=\scriptsize\sffamily, color=black!60, align=center] at (4.5, -0.8) {
    확률적 다양성 원리: 합집합(28) $>$ 개별 최대(22), 다양성 격차 = +27.3\%
};

\end{tikzpicture}
}
\caption{반복 횟수에 따른 성능 변화 (gpt-4o, Level~1). 개별 실행 성능은 반복 횟수에 따라 미미한 변동을 보이나, 누적 합집합은 꾸준히 증가한다. 이는 반복의 깊이보다 생성의 다양성이 더 효과적임을 시사한다.}
\label{fig:ablation}
\end{figure}

\begin{figure}[t]
\centering
\resizebox{\columnwidth}{!}{%
\begin{tikzpicture}
\begin{axis}[
    width=10cm,
    height=6.5cm,
    xlabel={반복 횟수 (Iterations)},
    ylabel={해결 문제 수 (/100)},
    xlabel style={font=\small\sffamily},
    ylabel style={font=\small\sffamily},
    xticklabel style={font=\small\sffamily},
    yticklabel style={font=\small\sffamily},
    xtick={1,3,5,10},
    ytick={0,5,10,15,19,20,21,22,25,28,30},
    ymin=0, ymax=32,
    xmin=0, xmax=11.5,
    grid=major,
    grid style={dashed, gray!25},
    legend style={
        at={(0.03,0.97)},
        anchor=north west,
        font=\small\sffamily,
        draw=gray!40,
        fill=white,
        fill opacity=0.9,
    },
]

% 개별 실행 성능
\addplot[
    name path=individual,
    color=blue!70,
    mark=*,
    mark size=3pt,
    thick,
    line join=round,
] coordinates {
    (1, 19) (3, 21) (5, 19) (10, 22)
};
\addlegendentry{개별 실행 성능}

% Union (누적 합집합) - 점선으로 표시
\addplot[
    name path=union,
    color=red!70,
    mark=triangle*,
    mark size=3.5pt,
    thick,
    dashed,
] coordinates {
    (1, 19) (3, 25) (5, 26) (10, 28)
};
\addlegendentry{누적 합집합 (Union)}

% 다양성 격차(shaded region)
\addplot[fill=green!10, draw=none, forget plot]
    fill between[of=union and individual];

% 개별 최대값 수평선
\addplot[
    color=blue!30,
    dashed,
    thin,
    domain=0:11.5,
] {22};

% Union 최대값 수평선
\addplot[
    color=red!30,
    dashed,
    thin,
    domain=0:11.5,
] {28};

% 차이 주석
\draw[{Stealth[length=4pt]}-{Stealth[length=4pt]}, thick, color=green!50!black]
    (10.8, 22) -- (10.8, 28);
\node[font=\scriptsize\bfseries\sffamily, color=green!50!black, right] at (10.8, 25) {+6};

% 라벨
\node[font=\tiny\sffamily, color=blue!50, anchor=west] at (11, 22) {max=22};
\node[font=\tiny\sffamily, color=red!50, anchor=west] at (11, 28) {union=28};

\end{axis}

% 하단 설명
\node[font=\scriptsize\sffamily, color=black!60, align=center] at (4.5, -0.8) {
    확률적 다양성 원리: 합집합(28) $>$ 개별 최대(22), 다양성 격차 = +27.3\%
};

\end{tikzpicture}
}
\caption{반복 횟수에 따른 성능 변화 (gpt-4o, Level~1). 개별 실행 성능은 반복 횟수에 따라 미미한 변동을 보이나, 누적 합집합은 꾸준히 증가한다. 이는 반복의 깊이보다 생성의 다양성이 더 효과적임을 시사한다.}
\label{fig:ablation}
\end{figure}
